\subsection{発展的な読み物}
\par\smallskip\noindent\refstepcounter{table}\label{tab:ch18-16}表 \thetable: 発展的な読み物の整理\par\smallskip
表 \thetable は、発展的な読み物の整理を比較・確認しやすい形で整理したものである。\par\smallskip
\begin{longtable}{>{\raggedright\arraybackslash}p{0.31\linewidth}>{\raggedright\arraybackslash}p{0.61\linewidth}}
\toprule
文献 & 読むと得られること \\
\midrule
\endfirsthead
\toprule
文献 & 読むと得られること \\
\midrule
\endhead
Abran, Software Metrics and Software Metrology & 測度、測定方法、測定結果という用語を厳密に使い分け、ソフトウェア測定をより正確に理解できる。測定の節を深めるのに有用である。 \\
Vincenti, What Engineers Know and How They Know It & 実際の工学事例を通じて、工学者が何を知り、どのように知識を得るかを理解できる。ソフトウェア以外の工学と比較するのに有用である。 \\
\bottomrule
\end{longtable}
